\usepackage{fontspec}
\usepackage{xeCJK}
\setmainfont{Latin Modern Roman}
\setCJKmainfont[ItalicFont={Noto Serif CJK SC},BoldItalicFont={Noto Serif CJK SC Bold}]{Noto Serif CJK SC}
\setCJKsansfont{Noto Sans CJK SC}
\setCJKmonofont{Noto Sans Mono CJK SC}
\newfontfamily\proseunicodefont{DejaVu Serif}
\usepackage[margin=1in]{geometry}
\usepackage{microtype}
\usepackage{amsmath,amssymb,amsthm,mathtools,bm}
\usepackage{graphicx}
\usepackage{xcolor}
\usepackage{booktabs}
\usepackage{array}
\usepackage{longtable}
\usepackage{enumitem}
\usepackage{float}
\usepackage{caption}
\usepackage{hyperref}
\usepackage{fancyhdr}
\usepackage{tikz}
\usetikzlibrary{arrows.meta,positioning,calc,decorations.pathmorphing}

\hypersetup{
  colorlinks=true,
  linkcolor=blue!50!black,
  urlcolor=blue!50!black,
  citecolor=blue!50!black
}

\setlength{\parskip}{0.5em}
\setlength{\parindent}{0pt}
\setlist[itemize]{topsep=0.25em,itemsep=0.2em}
\setlist[enumerate]{topsep=0.25em,itemsep=0.2em}
\setlength{\headheight}{40pt}
\numberwithin{equation}{chapter}

\newtheorem{theorem}{定理}[chapter]
\newtheorem{proposition}[theorem]{命题}
\newtheorem{lemma}[theorem]{引理}
\newtheorem{corollary}[theorem]{推论}
\theoremstyle{definition}
\newtheorem{definition}[theorem]{定义}
\newtheorem{example}[theorem]{例}
\theoremstyle{remark}
\newtheorem{remark}[theorem]{注}

\providecommand{\ket}[1]{\left\lvert#1\right\rangle}
\providecommand{\bra}[1]{\left\langle#1\right\rvert}
\providecommand{\braket}[2]{\left\langle#1\,\middle|\,#2\right\rangle}
\providecommand{\expect}[1]{\left\langle#1\right\rangle}
\newcommand{\lectureoverview}[1]{\paragraph{概述。} #1}
\newenvironment{classroomqa}{%
  \par\medskip
  \begin{quote}\small
}{%
  \end{quote}
  \medskip
}
\newcommand{\audiencequestion}[1]{\noindent\textbf{听众提问。} #1\par\smallskip}
\newcommand{\classquestion}[1]{\noindent\textbf{课堂提问。} #1\par\smallskip}
\newcommand{\lecturerresponse}[1]{\noindent\textbf{回答。} #1\par}
\newcommand{\lecturetimestamp}[1]{\footnote{讲课时间： #1.}}
\newcommand{\lectureframe}[1]{\par\smallskip{\centering\footnotesize\itshape 讲课画面： #1.\par}}

\makeatletter
\newcommand{\course@headerfont}{\normalfont\small\itshape}
\newcommand{\course@pagefont}{\normalfont\small}
\newcommand{\course@headerraise}{14pt}
\newlength{\course@headerboxheight}
\setlength{\course@headerboxheight}{\dimexpr\headheight-\course@headerraise-9pt\relax}
\newcommand{\course@leftheadbox}[1]{%
  \raisebox{\course@headerraise}{%
    \parbox[t][\course@headerboxheight][t]{0.82\textwidth}{%
      \raggedright\course@headerfont\strut\MakeUppercase{#1}\strut%
    }%
  }%
}
\newcommand{\course@rightheadbox}[1]{%
  \raisebox{\course@headerraise}{%
    \parbox[t][\course@headerboxheight][t]{0.82\textwidth}{%
      \raggedleft\course@headerfont\strut\MakeUppercase{#1}\strut%
    }%
  }%
}
\renewcommand{\chaptermark}[1]{\markboth{\MakeUppercase{\chaptername\ \thechapter.\ #1}}{}}
\renewcommand{\sectionmark}[1]{\markright{\MakeUppercase{\thesection\hspace{0.4em}#1}}}
\fancyhf{}
\fancyhead[L]{\course@leftheadbox{\leftmark}}
\fancyhead[R]{\course@pagefont\thepage}
\renewcommand{\headrulewidth}{0.5pt}
\pagestyle{fancy}
\fancypagestyle{plain}{%
  \fancyhf{}
  \fancyhead[L]{\course@leftheadbox{\leftmark}}
  \fancyhead[R]{\course@pagefont\thepage}
  \renewcommand{\headrulewidth}{0.5pt}
}
\makeatother

\renewcommand{\contentsname}{目录}
\renewcommand{\chaptername}{章}
\renewcommand{\figurename}{图}
\renewcommand{\tablename}{表}
\renewcommand{\proofname}{证明}
\renewcommand{\bibname}{参考文献}
\emergencystretch=2em
\pdfstringdefDisableCommands{\def\proseunicodefont{}}
